% =====================================================================
% Extinction dynamics: quasi-stationarity and first-passage times
% Insertable section for the Advanced Course in Continuous-Time
% Stochastic Processes. Requires amsmath and amssymb.
% =====================================================================
\section{Extinction dynamics: quasi-stationary distributions and first-passage times}
\label{sec:extinction-qsd-mfpt}

In a birth--death process, we can often calculate a stationary probability
 distribution by balancing the flows of probability between states. Yet this
approach raises a fundamental difficulty when a population cannot recover after
extinction. If immigration is absent, the empty state may be \emph{absorbing}:
once the last individual dies, the population remains extinct forever. What
can a stationary distribution tell us about such a population? And how long
will it take to become extinct?

We develop two complementary tools. The \emph{quasi-stationary distribution}
(QSD) describes the abundance of a population \emph{conditional on its still
being present}. The \emph{backward master equation} calculates the mean time
to reach a specified state, such as extinction, starting from a given initial
population size. Both descriptions originate from the same birth and death
transition rates.

\subsection{Absorbing states, survival, and extinction}
\label{subsec:absorbing-extinction}

Consider a continuous-time birth--death process $N(t)$ with state space
$n=0,1,2,\ldots$ and elementary transitions
\begin{equation}
 n\longrightarrow n+1\quad\text{at rate }g_n,\qquad
 n\longrightarrow n-1\quad\text{at rate }r_n.
 \label{eq:ext-transition-rates}
\end{equation}
The forward master equation for $n\geq 1$ is
\begin{equation}
 \frac{dP(n,t)}{dt}=
 g_{n-1}P(n-1,t)+r_{n+1}P(n+1,t)
 -(g_n+r_n)P(n,t),
 \label{eq:ext-forward}
\end{equation}
with the boundary equation
\begin{equation}
 \frac{dP(0,t)}{dt}=r_1P(1,t)-g_0P(0,t).
 \label{eq:ext-boundary}
\end{equation}
If $g_0=0$, the state $0$ is absorbing. More generally, immigration at
positive population sizes is compatible with $g_0=0$; what matters for
absorption is that no event can leave state $0$. In an absorbing process,
$P(0,t)$ cannot decrease, and the point mass concentrated at $n=0$ is
always a stationary distribution. We should therefore avoid saying that
``there is no stationary distribution'': the relevant question is whether a
\emph{nontrivial} stationary distribution on positive population sizes
exists. Extinction is certain in many familiar finite-population models,
but not for every infinite-state birth--death process; for example, a
supercritical branching process can escape extinction with positive
probability.

Let $\tau_0$ be the first time at which the population reaches zero:
\begin{equation}
 \tau_0=\inf\{t\geq0:N(t)=0\}.
 \label{eq:ext-hitting-def}
\end{equation}
Assume that the population starts in a positive state. Because zero is
absorbing, the \emph{survival probability} is
\begin{equation}
 S(t)=\Pr\{\tau_0>t\}=\Pr\{N(t)>0\}=1-P(0,t).
 \label{eq:ext-survival}
\end{equation}
From Eq.~\eqref{eq:ext-boundary}, for $g_0=0$ we obtain the exact relation
\begin{equation}
 \boxed{\frac{dS(t)}{dt}=-r_1P(1,t).}
 \label{eq:ext-survival-flux}
\end{equation}
Extinction can occur only through the last death, the jump from one
individual to none. This simple balance relation is the starting point
for both approaches developed below.

\subsection{Quasi-stationary distributions: abundance conditional on survival}
\label{subsec:ext-qsd}

Although $P(n,t)$ may ultimately concentrate at the absorbing state, the
abundance distribution among populations that have \emph{not yet} become
extinct may approach a stable shape. For $n\geq1$, define
\begin{equation}
 P_c(n,t)=\Pr\{N(t)=n\mid\tau_0>t\}
          =\frac{P(n,t)}{S(t)}.
 \label{eq:ext-conditional}
\end{equation}
The conditional probabilities satisfy $\sum_{n\geq1}P_c(n,t)=1$.
A \emph{quasi-stationary distribution} $\pi=(\pi_1,\pi_2,\ldots)$
is a normalized distribution over positive states that remains unchanged
when conditioned on survival. In particular, if the conditional distribution
converges as $t\to\infty$, its limit is a QSD:
\begin{equation}
 \pi_n=\lim_{t\to\infty}P_c(n,t),\qquad
 \sum_{n\geq1}\pi_n=1.
 \label{eq:ext-qsd-def}
\end{equation}
Such a limit need not exist in every process, and a QSD need not be unique.

We can derive an equation for the conditional probabilities without any new
assumptions. Differentiating Eq.~\eqref{eq:ext-conditional} gives
\begin{equation}
 \frac{dP_c(n,t)}{dt}
 =\frac{1}{S(t)}\frac{dP(n,t)}{dt}
  -\frac{P(n,t)}{S(t)^2}\frac{dS(t)}{dt}.
 \label{eq:ext-quotient}
\end{equation}
Using $dS/dt=-r_1P(1,t)=-r_1 S(t)P_c(1,t)$ yields
\begin{equation}
 \frac{dP_c(n,t)}{dt}
 =\frac{1}{S(t)}\frac{dP(n,t)}{dt}
  +r_1P_c(1,t)P_c(n,t).
 \label{eq:ext-conditional-general}
\end{equation}
Notice the \emph{positive} sign of the additional term. It compensates
for the probability continuously lost to extinction before the remaining
probabilities are renormalized.

For $n>1$, substitution of the forward master equation gives
\begin{align}
 \frac{dP_c(n,t)}{dt}={}&
 g_{n-1}P_c(n-1,t)+r_{n+1}P_c(n+1,t)\notag\\
 &-(g_n+r_n)P_c(n,t)+r_1P_c(1,t)P_c(n,t).
 \label{eq:ext-qsd-n}
\end{align}
For $n=1$, the incoming term from $n=0$ is absent because $g_0=0$:
\begin{equation}
 \frac{dP_c(1,t)}{dt}
 =r_2P_c(2,t)-(g_1+r_1)P_c(1,t)
  +r_1[P_c(1,t)]^2.
 \label{eq:ext-qsd-one}
\end{equation}
Setting the left-hand sides to zero gives the equations satisfied by a
QSD, whenever one exists.

\paragraph{An iterative construction.}
For simple one-dimensional processes, the conditional equations suggest
an iterative construction of the quasi-stationary distribution. Following
the procedure introduced in the original notes, we proceed as follows:
\begin{enumerate}
\item Make an initial guess for $P_c^*(1)=q$.
\item Set Eq.~\eqref{eq:ext-qsd-one} equal to zero and calculate $P_c^*(2)$:
\begin{equation}
 P_c^*(2)=\frac{(r_1+g_1)q-r_1q^2}{r_2}.
 \label{eq:ext-qsd-pi2}
\end{equation}
\item Set Eq.~\eqref{eq:ext-qsd-n} equal to zero and successively calculate
$P_c^*(n)$, for $n>2$, using
\begin{equation}
 P_c^*(n)=\frac{(r_{n-1}+g_{n-1}-r_1q)P_c^*(n-1)
                   -g_{n-2}P_c^*(n-2)}{r_n}.
 \label{eq:ext-qsd-recursion}
\end{equation}
Continue the recurrence until the probabilities become sufficiently
small, according to a chosen numerical tolerance $\varepsilon$.
\item Normalize the resulting positive-state probabilities and obtain an
updated estimate of $P_c^*(1)$. Compare it with the initial guess $q$;
if they differ appreciably, repeat the procedure with the new estimate.
\end{enumerate}

This four-step procedure conveys the constructive idea behind the
quasi-stationary solution. In numerical applications, however, stopping
when one probability becomes small does not by itself control the
remaining tail, and normalization does not guarantee convergence or
nonnegative probabilities. The iteration should therefore be checked
for stability and positivity; for finite truncated state spaces, an
eigenvector calculation is often a more reliable alternative.

\subsection{Extinction times from quasi-stationarity}
\label{subsec:ext-qsd-time}

We now ask a natural question: \emph{how long does it take, on average,
for a population to become extinct?}. It is
helpful to begin with $P(0,t)$ rather than directly introducing a first-passage calculation. Assume that $N(0)>0$ and that $g_0=0$, so that zero is absorbing. A population is at zero at time $t$ if and only if it has already become extinct by that time. Thus $P(0,t)$ is precisely the cumulative distribution function of the random variable
$\tau_0$, the \emph{time to extinction}:
\begin{equation}
 \Pr\{\tau_0\leq t\}=P(0,t),
 \qquad \Pr\{\tau_0>t\}=1-P(0,t)=S(t).
 \label{eq:ext-cdf-time}
\end{equation}
When the extinction-time distribution has a density, it is given by
$f_{\tau_0}(t)=dP(0,t)/dt$. Consequently, if extinction occurs almost
surely and the mean is finite, we can write
\begin{equation}
 \langle \tau_0\rangle
 =\int_0^\infty t\,\frac{dP(0,t)}{dt}\,dt.
 \label{eq:ext-mean-cdf-density}
\end{equation}
Integrating by parts over a finite interval $[0,T]$ gives
\begin{align}
 \int_0^T t\,\frac{dP(0,t)}{dt}\,dt
 &=\int_0^T t\,[-S'(t)]\,dt\\
 &=-T S(T)+\int_0^T S(t)\,dt.
 \label{eq:ext-parts-finite}
\end{align}
For a finite mean extinction time, $T S(T)\to0$ as $T\to\infty$.
We therefore obtain the particularly useful identity
\begin{equation}
 \boxed{\langle \tau_0\rangle
   =\int_0^\infty [1-P(0,t)]\,dt
   =\int_0^\infty S(t)\,dt.}
 \label{eq:ext-survival-integral}
\end{equation}
This is the geometric interpretation of the mean extinction time:
it is the area under the survival-probability curve. More generally,
the survival-integral identity remains valid as an extended-valued
expectation even when the mean is infinite (and also if extinction
is not certain). The derivation using a probability density above
is particularly transparent in the usual case of almost-sure
extinction with finite mean.
\smallskip

We can now approximate this area using the quasi-stationary
probabilities obtained in the previous subsection. Combining the
boundary master equation with $P(1,t)=P_c(1,t)[1-P(0,t)]$, we find
\begin{equation}
 \frac{dP(0,t)}{dt}
 =r_1P(1,t)
 =r_1P_c(1,t)[1-P(0,t)].
 \label{eq:ext-hazard-p0}
\end{equation}
Equivalently, the survival probability obeys
\begin{equation}
 \boxed{\frac{dS(t)}{dt}=-r_1P_c(1,t)S(t).}
 \label{eq:ext-hazard}
\end{equation}
The quantity $r_1P_c(1,t)$ is the instantaneous extinction hazard,
conditional on survival: extinction requires first having a single
individual and then losing that last individual.

If the initial population size is drawn from a quasi-stationary
distribution $\pi$, then $P_c(1,t)=\pi_1$ for every $t\geq0$.
Since $S(0)=1$, integration yields
\begin{equation}
 1-P(0,t)=S(t)=\exp(-r_1\pi_1t).
 \label{eq:ext-qsd-survival}
\end{equation}
Substitution into Eq.~\eqref{eq:ext-survival-integral} gives
\begin{equation}
 \boxed{\langle \tau_0\rangle_{\pi}
    =\int_0^\infty\exp(-r_1\pi_1t)\,dt
    =\frac{1}{r_1\pi_1}.}
 \label{eq:ext-qsd-mean}
\end{equation}
This expression is \emph{exact} when the process starts in the QSD.
For a population starting from a specified positive abundance,
replacing $P_c(1,t)$ by $\pi_1$ gives an \emph{approximation} whose
accuracy depends on how rapidly the conditional distribution approaches
quasi-stationarity compared with the characteristic extinction time.
The initial transient may be important for some starting abundances.

This forward-equation argument begins with the entire distribution
$P(n,t)$ and uses its probability flux into the absorbing state.
We next develop a complementary approach: determining the mean time
to extinction directly from the initial abundance, by means of the
\emph{backward master equation}.

\subsection{Backward master equations and mean first-passage times}
\label{subsec:ext-backward}

The forward master equation describes the probability of occupying state
$n$ at time $t$. A different question is: \emph{starting from a given
state $n$, how long will it take to reach a particular target for the
first time?} The answer is the \emph{mean first-passage time} (MFPT),
\begin{equation}
 T_n=\mathbb E[\tau_A\mid N(0)=n],\qquad
 \tau_A=\inf\{t\geq0:N(t)\in A\},
 \label{eq:ext-mfpt-def}
\end{equation}
for a target set $A$. For extinction, $A=\{0\}$ and $T_0=0$.

\subsubsection{First-event argument: the general backward equation}
For $n\notin A$, the total rate at which an event occurs is
$R_n=g_n+r_n$. Assuming $R_n>0$, the time until the next event has
mean $1/R_n$. The next event increases the population with probability
$g_n/R_n$ and decreases it with probability $r_n/R_n$.
Conditioning on this first event gives
\begin{equation}
 T_n=\frac{1}{R_n}
   +\frac{g_n}{R_n}T_{n+1}
   +\frac{r_n}{R_n}T_{n-1}.
 \label{eq:ext-first-step}
\end{equation}
Rearranging,
\begin{equation}
 \boxed{-1=g_n(T_{n+1}-T_n)+r_n(T_{n-1}-T_n),
       \qquad n\notin A.}
 \label{eq:ext-backward-equation}
\end{equation}
The boundary condition is $T_n=0$ for $n\in A$. A finite state space may
require an additional upper-boundary condition; for an infinite state
space, further conditions are sometimes needed to select a finite
solution. The first-passage mean may also be infinite.

Why is the equation called \emph{backward}? The unknown quantity $T_n$
is indexed by the \emph{initial} state of the process, whereas the
forward equation follows the probability of the state \emph{reached}
after a given time. Both are consequences of the same transition rates.

\subsubsection{Worked example: extinction in a pure-death process}
Suppose each individual dies independently at per-capita rate $\mu$,
there are no births, and there is no immigration:
\begin{equation}
 g_n=0,\qquad r_n=\mu n,\qquad T_0=0.
 \label{eq:ext-pure-death}
\end{equation}
The backward equation reduces to
\begin{equation}
 -1=\mu n(T_{n-1}-T_n),\qquad
 T_n=\frac{1}{\mu n}+T_{n-1}.
 \label{eq:ext-pure-death-step}
\end{equation}
Applying this recurrence repeatedly gives
\begin{equation}
 \boxed{T_n=\frac{1}{\mu}\sum_{k=1}^n\frac{1}{k}
             =\frac{H_n}{\mu}.}
 \label{eq:ext-pure-death-result}
\end{equation}
Here $H_n$ denotes the $n$th harmonic number. For large $n$,
$H_n\simeq\ln n+\gamma$, where $\gamma$ is Euler's constant. Thus the
mean extinction time grows roughly logarithmically with initial
population size. In this simple case the backward method gives an
exact extinction time without solving the time-dependent forward
master equation.

\subsubsection{Worked example: time to reach two individuals}
Consider a site with immigration rate $\lambda>0$, per-capita birth
rate $\beta\geq0$, and per-capita death rate $\mu\geq0$. How long,
on average, until a population starting at $n=0$ first reaches $n=2$?
Only states 0 and 1 need to be considered before the target is reached.
The rates are $g_0=\lambda$, $g_1=\lambda+\beta$,
$r_0=0$, $r_1=\mu$; set $T_2=0$. The backward equations read
\begin{align}
 -1&=\lambda(T_1-T_0),\label{eq:ext-two-zero}\\
 -1&=(\lambda+\beta)(T_2-T_1)+\mu(T_0-T_1).
 \label{eq:ext-two-one}
\end{align}
Solving,
\begin{equation}
 \boxed{T_1=\frac{1+\mu/\lambda}{\lambda+\beta},\qquad
 T_0=\frac{1}{\lambda}+\frac{1+\mu/\lambda}{\lambda+\beta}.}
 \label{eq:ext-two-solution}
\end{equation}
The target $n=2$ need not be an absorbing state of the original process:
we simply stop the clock at the first arrival there. This example
shows why first-passage methods apply to recovery and colonization as
well as extinction.

\subsection{Applications and comparison of approaches}
\label{subsec:ext-applications}

\paragraph{Application: population collapse under harvesting.}
In an exploited resource population, the transition rates may combine
reproduction, density-dependent mortality, immigration, and saturating
harvesting. One illustrative choice is
\begin{equation}
 g_n=\lambda+bn,\qquad
 r_n=dn+\frac{b}{K}n(n-1)+\frac{hn}{a+n}.
 \label{eq:ext-harvest-example}
\end{equation}
Here the last term represents Holling-type-II harvesting. These rates
are presented only as an \emph{illustrative model}, not as the exact
microscopic rates in any particular publication. For collapse, choose a
low-abundance threshold $n_c$ and set $T_n=0$ for $n\leq n_c$; for
recovery, select an upper threshold instead. The backward equation
then gives a set of linear relations for the expected threshold-crossing
times, as in the work of Cuenda, Llorente, and Capit\'an (2020).
If $\lambda>0$, zero is no longer absorbing: the distinction between
\emph{first arrival at zero} and \emph{permanent extinction} must then
be kept explicit.

\paragraph{How the two approaches complement one another.}
The QSD characterizes \emph{how many individuals are typically present,
given that extinction has not occurred}. Once this conditional
distribution is stationary, the extinction hazard is constant and its
inverse gives the exact mean extinction time for QSD initial data.
The backward equation instead asks \emph{how long it takes to reach a
target from a specified initial state}. It need not invoke
quasi-stationarity and applies to any reachable target, not just zero.

As a useful synthesis, suppose the state $0$ is absorbing and the
process has a QSD $\pi$. The MFPT $T_n$ from each initial state, if
finite, satisfies Eq.~\eqref{eq:ext-backward-equation}. Averaging these
mean times over initial states drawn from $\pi$ must recover the QSD
result:
\begin{equation}
 \boxed{\sum_{n\geq1}\pi_nT_n=\frac{1}{r_1\pi_1},}
 \label{eq:ext-bridge}
\end{equation}
provided these quantities are finite. This equality links the forward
conditional distribution and the backward first-passage analysis.

\subsubsection*{Exercises}
\begin{enumerate}
 \item \textbf{Absorbing boundary and survival.} Starting from the
 boundary master equation at $n=0$, derive
 $dS/dt=-r_1P(1,t)$ when $g_0=0$. What changes if immigrants can
 arrive at an empty site ($g_0>0$)?

 \item \textbf{Conditional probability balance.} Differentiate
 $P_c(n,t)=P(n,t)/S(t)$ and recover Eqs.~\eqref{eq:ext-qsd-n}
 and \eqref{eq:ext-qsd-one}. Explain in words why the conditional
 master equation contains a positive normalization term.

 \item \textbf{Quasi-stationarity and extinction.} Assume that
 $P_c(n,t)=\pi_n$ for all $t\geq0$. Derive the survival function and
 show that the mean extinction time is $1/(r_1\pi_1)$. Why is this
 result not necessarily exact for a population starting from a fixed
 size $n_0$?

 \item \textbf{Pure-death extinction time.} Derive $T_1$, $T_2$,
 and $T_3$ from the backward recurrence. Check them against
 $T_n=H_n/\mu$. Why do additional initial individuals contribute
 progressively smaller increments to the mean extinction time?

 \item \textbf{A threshold-crossing problem.} Consider immigration,
 birth, and death with rates $g_n=\lambda+\beta n$ and $r_n=\mu n$.
 Derive the backward equations for the first passage from $n=0$
 to $n=2$, and verify Eq.~\eqref{eq:ext-two-solution}. Discuss what
 happens when $\mu$ increases.

 \item \textbf{Optional numerical project.} Choose a finite-capacity
 birth--death process with $g_0=0$ and extinction certain. Estimate
 its QSD from the conditional forward equations or the transient
 generator, and calculate the MFPTs from the backward equations.
 Compare $\sum_{n\geq1}\pi_nT_n$ with $1/(r_1\pi_1)$.
\end{enumerate}

\medskip
\noindent\textbf{Take-home message.} The same elementary birth and death
rates support two distinct but complementary questions: \emph{what does
a surviving population look like?} and \emph{how long does it take to
reach extinction or another threshold?} Quasi-stationarity addresses
the first; backward first-passage equations address the second.
